Deux flûtistes jouent la même note sur des flûtes identiques. La flûte de Léa est à la température de la salle, \qty{20}{\degreeCelsius}, et produit un son de \qty{\v{f}}{Hz}. La flûte de Samuel, qu'il vient de réchauffer en soufflant longtemps dedans, contient de l'air à \qty{\v{T2}}{\degreeCelsius}. On utilise $v = \qty{331}{m/s}\sqrt{T/\qty{273}{K}}$.

\begin{enumerate}[label=\alph*)]
    \item Expliquez pourquoi la fréquence de résonance d'une flûte est proportionnelle à $\sqrt{T}$, puis déterminez la fréquence produite par la flûte de Samuel.
    \item Quelle est la fréquence des battements entendus quand ils jouent ensemble ?
    \item Samuel doit-il allonger ou raccourcir sa flûte (en tirant sur la tête de l'instrument) pour éliminer les battements ? Justifiez.
\end{enumerate}
